\section{模型的评价}

\subsection{模型的优点}
\begin{itemize}[itemindent=2em]
\item \textbf{口径统一：} 四个问题共用同一条能源结算模型与同一可行任务族，结论可直接跨问题比较，避免了"换口径换结论"。
\item \textbf{可验证：} 解析与 LP/MILP 逐值一致、SOC 守恒到机器精度、8 类约束审计全通过，结果可复现，附完整源程序与支撑材料。
\item \textbf{诚实与可解释：} 明确报告 ABC 储能零价值、碳约束退化、任务储能替代关系等"反直觉但真实"的结论，并用边际成本面给出物理可解释。
\item \textbf{稳健性显式刻画：} 通过 $\rho$ 扫描定位任务迁移收益的转折区间，给出两类灵活性的适用边界。
\end{itemize}

\subsection{模型的缺点}
\begin{itemize}[itemindent=2em]
\item \textbf{任务侧启发式：} F10/F11 采用冻结边际成本引导的贪心改进，未提供全局最优下界，交互项 $I$ 为点估计而非认证区间。
\item \textbf{预测信息量有限：} 到达序列噪声大，岭回归仅略优于均值基准，预测对未来负荷的支撑价值有限。
\item \textbf{线性近似：} 储能与功率映射做了线性化/常数效率处理，未考虑充放电损耗随功率的非线性与电池寿命衰减。
\end{itemize}
